\textbf{Задача 5}  Язык называется унарным, если все слова в нём имеют вид $1^k$. Докажите, что если существует унарный $\np$ полный язык, то $\p=\np$.

\begin{proof} 
Пусть этот унарный $\np$ полный язык $A$.
Так как он $\np$ полный, то к нему сводится $\sat$, то есть существует машина $g(x)$ работающая за время $\le p(|x|)$, такая что
\[\forall \varphi\quad \varphi\in\sat \iff g(\varphi)\in A \tag{1}\]

\begin{statement}   Если для формул $\varphi_1, \varphi_2$ имеем $g(\varphi_1)= g(\varphi_2)$, то 
\[\varphi_1 \vee \varphi_2 \in \sat \iff \varphi_1 \in \sat \iff \varphi_2 \in \sat\]
\end{statement}
\begin{proof}
    Вторая равносильность получается из (1) так:
    \[\varphi_1\in\sat \iff g(\varphi_1)\in A \iff g(\varphi_2)\in A \iff \varphi_2\in\sat \]
    Первая равносильность, следует из того что если формула верна, то верна либо $\varphi_2$ или $\varphi_1$, и по второй равносильности в любом случае следует, что $\varphi_1$ выполнима. В обратную сторону очевидно.
\end{proof}

Покажем, что $\sat\in\p$. Зафиксируем формулу $\varphi(x_1,\ldots, x_m),\ n=|\varphi|, m \le n$.

Будем по очереди фиксировать значения $x_i$. Пусть $\psi_0=\psi'_0=\varphi$. Положим 
\[\psi'_{k+1} = \psi_k|_{x_{k+1}=0} \vee \psi_k|_{x_{k+1}=1}\]

Используя предыдущее утверждение, можно сократить количество подформул  (подформула --- $\varphi$ с некотрыми фиксированными переменными) в $\psi'_{k + 1}$. Так как у нас не более $2 \cdot p(n)$ различных значений $g$ (считаем не унарные значения за одно), то кол-во подформул в $\psi'_{k + 1}$ можно сократить до не более чем $2 \cdot p(n)$ подформул. Обозначим это сокращение как $\psi_{k+1}$

Имеем 
\[\varphi\in\sat\iff\psi'_k\in\sat\iff\psi_k\in\sat \quad \forall k\]

Но в $\psi_m$ нету переменных, поэтому \[\varphi\in\sat\iff\psi_m\in\sat\iff\psi_m=1\]
Последнее можно просто посчитать.

Оценим время работы.

Количество подформул в $\psi_k$ и в $\psi'_k$ не более, чем $C\cdot p(n)$
Длина каждой подформулы $\le n$

всего $\psi$ $O(m)=O(n)$ штук. Пересчет также работает за полином от $n$. Значит итоговое время работы будет полиномиально от $n$. Значит $\sat\in\p$. 
А так как $\sat\ \np$ полный, то $\p=\np$ 
\end{proof}